Small tabular datasets are a regime where practitioners regularly reach for L2 regularization, but the size and consistency of its benefit are often unclear. We study a simple, controlled version of that question: does validation-tuned L2 regularization improve held-out generalization relative to an otherwise matched logistic-regression baseline on datasets with fewer than 1000 samples? We evaluate four real classification datasets with stratified $70/15/15$ train/validation/test splits and three fixed random seeds per dataset. For each split, we compare a near-unregularized logistic regression baseline, a validation-tuned L2 logistic regression model, and a stratified dummy classifier. We report test accuracy, weighted F1, train-validation and train-test gaps, coefficient norms, and paired statistical tests over the 12 dataset-seed pairs.

L2 regularization improves the mean test accuracy from $0.8292$ to $0.8385$, but that gain is not statistically significant ($p=0.345$). The same pattern holds for weighted F1 ($0.8276$ to $0.8357$, $p=0.431$). In contrast, L2 sharply reduces the mean train-validation gap from $0.1297$ to $0.0257$, a paired reduction of $0.1040$ with a $95\%$ confidence interval of $[-0.1792,-0.0287]$ and $p=0.011$. This gap shrinks in $11/12$ runs, while test accuracy improves in only $5/12$ runs.

Our main conclusion is practical rather than universal: on small real tabular datasets, L2 regularization is a dependable anti-overfitting tool, but not a reliable accuracy booster. That distinction matters when regularization is used to justify stronger claims about predictive performance.
